\documentclass{exam}
\usepackage{../../shah350}

\hypersetup{colorlinks=true, linktoc=section, linkcolor=blue}

\pagestyle{headandfoot}
\firstpageheadrule
\runningheadrule
\firstpageheader{Prof. Lepowsky \\ Linear Algebra}{Homework \#7}{\href{https://github.com/jeevanshah07/MATH350/blob/main/homework/homework7/main.tex}{\underline{Source}} \\ Jeevan Shah}
\runningheader{Linear Algebra \\ Homework \#7}{}{Shah}
\firstpagefooter{}{\thepage}{}
\runningfooter{ }{\thepage}{ }

\printanswers

% 4.1.9, 4.1.11a. 4.1.11*
% 4.2.7, 4.2.8, 4.2.14, 4.2.18, 4.2.23
% 4.4.2c, 4.4.3c, 4.4.3e

\begin{document}
  \colorbox{red}{\underline{$4.1.9$}:}
  \begin{solution}
    \begin{proof}
      Let $A, B \in \MM_{2\times 2}(\FF)$ with 
      \[
        A = \begin{pmatrix}
          a & b \\ c & d 
        \end{pmatrix}
        \mand B = \begin{pmatrix}
          e & f \\ g & h
        \end{pmatrix}
      \]
      so 
      \[
        AB = \begin{pmatrix}
          ae + bg & af + bh \\
          ce + dg & cf + dh
        \end{pmatrix}.
      \]
      Then, 
      \begin{align*}
        \det(AB) &= (ae + bg)(cf + dh) - (af + bh)(ce + dg) \\
        &= acef + adeh + bcfg + bdgh - (acef + adfg + bceh + bdgh) \\
        &= adeh - adfg + bcfg - bceh \\
        &= ad(eh - fg) - bc(eh - fg) \\
        &= (ad - bc)(eh - fg) \\
        &= \det(A) \cdot \det(B)
      \end{align*}
      which was to be shown.
    \end{proof}
  \end{solution}

  \colorbox{red}{\underline{$4.1.11a$}:}
  \begin{solution}
    \begin{proof}
      Consider some $E \in \MM_{2 \times 2}(\FF)$ an elementary matrix given by applying at most one elementary row operation to $I$, the identity matrix. 
      \newline\underline{Case 1 (no operation):} If no operations are applied then clearly $E = I$ so 
      \[
        \delta(E) = \delta(I) = 1 = \det(I) = \det(E).
      \]
      \newline\underline{Case 2 ($r_1 \leftrightarrow r_2$):} If the only operation applied is switching rows then 
      \[
        I = \begin{pmatrix}
          1 & 0 \\ 0 & 1
        \end{pmatrix} 
        \xrightarrow{r_1 \leftrightarrow r_2} 
        \begin{pmatrix}
          0 & 1 \\
          1 & 0 
        \end{pmatrix} = E
      \]
      so $\det(E) = -1$. Now, using the facts that $\delta$ is a linear function of each row and that $\delta(A) = 0$ if any two rows of $A$ are identical, we may consider 
      \begin{align*}
        0 = \delta\Biggl(\begin{pmatrix}
          1 & 1 \\ 1 & 1
        \end{pmatrix}\Biggr) 
        &= \delta\Biggl(\begin{pmatrix}
          1 & 0 \\
          1 & 1 
         \end{pmatrix}\Biggr)
         + \delta\Biggl(\begin{pmatrix}
          0 & 1 \\
          1 & 1 
         \end{pmatrix}\Biggr) \\
        &= \delta\Biggl(\begin{pmatrix}
          1 & 0 \\
          0 & 1 
         \end{pmatrix}\Biggr)
         + \delta\Biggl(\begin{pmatrix}
          1 & 0 \\
          1 & 0 
         \end{pmatrix}\Biggr)
         + \delta\Biggl(\begin{pmatrix}
          0 & 1 \\
          1 & 0 
         \end{pmatrix}\Biggr)
         + \delta\Biggl(\begin{pmatrix}
          0 & 1 \\
          0 & 1 
         \end{pmatrix}\Biggr) \\
        &= \delta(I) + 0 + \delta(E) + 0 = 0 \\
        &\Rightarrow\, \delta(E) = -1 = \det(E)
      \end{align*}
      \newline\underline{Case 3 ($r_i \to k r_i$):} In the case that a row is scaled, without loss of generality (since the case in which $i = 2$ is identical), we take $i = 1$ so $I \to E$ by $r_1 \to k r_1$ for $k \in \FF$: 
      \[
        I = \begin{pmatrix}
          1 & 0 \\
          0 & 1
        \end{pmatrix} 
        \xrightarrow{r_1 \to k r_1} 
        \begin{pmatrix}
          k & 0 \\
          0 & 1
        \end{pmatrix} = E
      \]
      so $\det(E) = k$. Then, 
      \[
        \delta\Biggl(\begin{pmatrix} 
          k & 0 \\ 0 & 1 
        \end{pmatrix}\Biggr)
        = \delta\Biggl(\begin{pmatrix} 
          k\cdot 1 & k\cdot 0 \\ 0 & 1 
        \end{pmatrix}\Biggr)
        = k\,\delta\Biggl(\begin{pmatrix} 
          1 & 0 \\ 0 & 1 
        \end{pmatrix}\Biggr)
        = k\,\delta(I) = k = \det(E).
      \]
      \newline\underline{Case 4 ($r_i \to r_i + k r_j$):} As before, without loss of generality we may take $i = 1$ and $j = 2$, as the argument is identical should they be swapped, and $k \in \FF$ so $I \to E$ by $r_1 \to r_1 + k r_2$: 
      \[
        I = \begin{pmatrix}
          1 & 0 \\
          0 & 1
        \end{pmatrix}
        \xrightarrow{r_1 \to r_1 + k r_2}
        \begin{pmatrix}
          1 & k \\
          0 & 1
        \end{pmatrix} = E
      \]
      so $\det(E) = 1$. But 
      \[
        \delta\Biggl(\begin{pmatrix}
          1 & k \\
          0 & 1
        \end{pmatrix}\Biggr) 
        = \delta\Biggl(\begin{pmatrix}
          1 & 0 \\
          0 & 1
        \end{pmatrix}\Biggr) 
        + \delta\Biggl(\begin{pmatrix}
          0 & k \\
          0 & 1
        \end{pmatrix}\Biggr)
        = \delta(I) + k\delta\Biggl(\begin{pmatrix}
          0 & 1 \\
          0 & 1
        \end{pmatrix}\Biggr)
        = 1 + 0 = 1 = \det(E).
      \]
      Hence, $\det(E) = \delta(E)$ for all elementary $2 \times 2$ matrices.
    \end{proof}
  \end{solution}

  \colorbox{red}{\underline{$4.1.11*$}:}
  \begin{solution}
    \begin{proof}
      Let $A \in \MM_{2 \times 2}(\FF)$ so 
      \[
        A = \begin{pmatrix}
          a & b \\ c & d
        \end{pmatrix}.
      \]
      So, 
      \begin{align*}
        \delta\Biggl(\begin{pmatrix}
          a & b \\
          c & d
        \end{pmatrix}\Biggr) &= \delta\Biggl(\begin{pmatrix}
          a & 0 \\
          c & d
        \end{pmatrix}\Biggr) + \delta\Biggl(\begin{pmatrix}
        0 & b \\
        c & d 
        \end{pmatrix}\Biggr) \\
        &= \delta\Biggl(\begin{pmatrix}
        a & 0 \\
        c & 0 
        \end{pmatrix}\Biggr) + \delta\Biggl(\begin{pmatrix}
        a & 0 \\
        0 & d 
        \end{pmatrix}\Biggr) + \delta\Biggl(\begin{pmatrix}
        0 & b \\
        c & 0 
        \end{pmatrix}\Biggr) + \delta\Biggl(\begin{pmatrix}
        0 & b \\
        0 & d 
        \end{pmatrix}\Biggr) \\
        &= 0 + ad + (-bc) + 0 \\
        &= ad - bc = \det(A)
      \end{align*}
      with the third step following from the second using identical techniques to the previous cases. \\
      Thus, $\delta(A) = \det(A)$ for all $A \in \MM_{2 \times 2}(\FF)$.
    \end{proof}
  \end{solution}

  \newpage

  \colorbox{red}{\underline{$4.2.7$}:}
  \begin{solution}
    \[
      \begin{vmatrix}
        0 & 1 & 2 \\
        -1 & 0 & -3 \\
        2 & 3 & 0 
      \end{vmatrix}
      = \begin{vmatrix}
        1 & 2 \\
        3 & 0 
      \end{vmatrix}
      + 3\begin{vmatrix}
        0 & 1 \\
        2 & 3
      \end{vmatrix}
      = -6 + 3(-2) = \boxed{-12}
    \]
  \end{solution}

  \colorbox{red}{\underline{$4.2.8$}:}
  \begin{solution}
    \[
      \begin{vmatrix}
        1 & 0 & 2 \\
        0 & 1 & 5 \\
        -1 & 3 & 0 
      \end{vmatrix}
      = -1\begin{vmatrix}
        0 & 2 \\
        1 & 5
      \end{vmatrix}
      -3 \begin{vmatrix}
        1 & 2 \\
        0 & 5
      \end{vmatrix}
      = -1(-2) - 3(5) = \boxed{-13}
    \]
  \end{solution}

  \colorbox{red}{\underline{$4.2.14$}:}
  \begin{solution}
    \[
      \begin{vmatrix}
        2 & 3 & 4 \\
        5 & 6 & 0 \\
        7 & 0 & 0
      \end{vmatrix}
      = 7\begin{vmatrix}
        3 & 4 \\
        6 & 0 
      \end{vmatrix}
      = 7(-24) = \boxed{-168}
    \]
  \end{solution}

  \colorbox{red}{\underline{$4.2.18$}:}
  \begin{solution}
    Consider the sequence of elementary row operations 
    \[
      A = \begin{pmatrix}
        1 & -2 & 3 \\
        -1 & 2 & -5 \\
        3 & -1 & 2
      \end{pmatrix}
      \xrightarrow{r_1 \to r_1 + r_2}
      \begin{pmatrix}
        0 & 0 & -2 \\
        -1 & 2 & -5 \\
        3 & -1 & 2
      \end{pmatrix}
      = \tilde{A}
    \]
    so $\det(A) = \det(\tilde{A})$. Then,
    \[
      \det(A) = \det(\tilde{A}) = \begin{vmatrix}
        0 & 0 & -2 \\
        -1 & 2 & -5 \\
        3 & -1 & 2
      \end{vmatrix} 
      = -2\begin{vmatrix}
        -1 & 2 \\ 3 & -1
      \end{vmatrix}
      = -2(-5) = \boxed{10}
    \]
  \end{solution}

  \colorbox{red}{\underline{$4.2.23$}:}
  \begin{solution}
    \begin{proof}
      We will prove this via induction on $n$ for $A \in \MM_{n \times n}(\FF)$. Consider the first non-trivial value of $n$, $n = 2$ and let $A \in \MM_{n \times n}(\FF)$ be any square upper-triangular matrix: 
      \[
        A = \begin{pmatrix}
          a & b \\
          0 & c
        \end{pmatrix}
      \]
      so $\det(A) = ac$ is the product of its diagonal entries. Thus, for $n = 2$ our proposition is true and we can now consider $n = 2$ as the base case. Now, suppose that for $n = k \geq 2$ the determinant of a $k \times k$ upper-triangular matrix is the product of its diagonal entries. We must now show this is true for $n = k+1$. Let $A \in \MM_{(k+1) \times (k+1)}(\FF)$ and $\tilde{A} \in \MM_{k \times k}(\FF)$ so that 
      \[
        A = \begin{pmatrix}
          a_{11} & \cdots & a_{1(k+1)} \\
          \vdots & \ddots & \vdots \\
          0 & \cdots & a_{(k+1)(k+1)}
        \end{pmatrix}
        \mand 
        \tilde{A} = \begin{pmatrix}
          a_{11} & \cdots & a_{1k} \\
          \vdots & \ddots & \vdots \\
          0 & \cdots & a_{kk}
        \end{pmatrix}
      \]
      where $A, \tilde{A}$ are both upper triangular. Now, we consider the determinant of $A$ found by cofactor expansion across the last row (the cofactor sign is $(-1)^{2(k+1)} = 1$) 
      \[
        \det(A) = a_{(k+1)(k+1)}\begin{vmatrix}
          a_{11} & \cdots & a_{1k} \\
          \vdots & \ddots & \vdots \\
          0 & \cdots & a_{kk}
        \end{vmatrix}
        = a_{(k+1)(k+1)}\det(\tilde{A}) = a_{(k+1)(k+1)}a_{kk} \cdots a_{22}a_{11}
      \]
      by the induction hypothesis. Thus, the determinant of any upper triangular square matrix is the product of its diagonal entries.
    \end{proof}
  \end{solution}


  \colorbox{red}{\underline{$4.3.12$}:}
  \begin{solution}
    \begin{proof}
      Let $Q \in \MM_{n \times n}(\FF)$ be orthogonal such that $QQ^{T} = I$. Then, 
      \begin{align*}
        \det(I) = \det(QQ^{T}) &= \det(Q) \cdot \det(Q^{T}) \tag{Thm~4.7} \\
        &= \det(Q) \cdot \det(Q) \tag{Thm~4.8} \\
        &= \left[\det(Q)\right]^{2} = 1 \\
        &\Rightarrow (\det(Q) - 1)(\det(Q) + 1) = 0
      \end{align*}
      and hence either $\det(Q) = 1$ or $\det(Q) = -1$.
    \end{proof}
  \end{solution}

  \colorbox{red}{\underline{$4.3.15$}:}
  \begin{solution}
    \begin{proof}
      Let $A, B \in \MM_{n \times n}(\FF)$ such that $A$ is similar to $B$. That is, there exists some invertible $P \in \MM_{n \times n}(\FF)$ such that $B = P^{-1}AP$. But 
      \begin{align*}
        \det(B) = \det(P^{-1}AP) &= \det(P^{-1}) \cdot \det(AP) \tag{Thm~4.7} \\
        &= \det(P^{-1}) \cdot \det(A) \cdot \det(P) \tag{Thm~4.7} \\
        &= \frac{1}{\det(P)} \cdot \det(A) \cdot \det(P) \tag{Corr. to Thm~4.8} \\
        &= \det(A)
      \end{align*}
      thus $\det(B) = \det(A)$ as was to be shown.
    \end{proof}
  \end{solution}

  \newpage 

  \colorbox{red}{\underline{$4.3.24$}:}
  \begin{solution}
    \begin{proof}
      Let $A$ be the matrix with $-1$ on the subdiagonal, last column $(a_0, \dots, a_{n-1})^T$, and zeros elsewhere, and consider $\tilde{A} = A + tI$ so 
      \[
        \tilde{A} = A + tI = \begin{pmatrix}
          t & 0 & 0 & \cdots & 0 & a_0 \\
          -1 & t & 0 & \cdots & 0 & a_1 \\
          0 & -1 & t & \cdots & 0 & a_2 \\
          \vdots & \vdots & \vdots & & \vdots & \vdots \\
          0 & 0 & 0 & \cdots & -1 & a_{n-1} + t
        \end{pmatrix}.
      \]
      Now, we will prove that 
      \[
        \det(\tilde{A}) = t^n + t^{n-1}a_{n-1} + \cdots + t^{2}a_2 + t a_1 + a_0
      \]
      via induction on $n$ for all matrices $\tilde{A} \in \MM_{n \times n}(\FF)$ such that $\tilde{A}$ is of the necessary form. Consider $n = 2$ as our first non-trivial value of $n$ and hence our base case. Then, 
      \[
        \tilde{A} = \begin{pmatrix}
          0 & a_0 \\
          -1 & a_1
        \end{pmatrix}
        + \begin{pmatrix}
          t & 0 \\
          0 & t
        \end{pmatrix}
        = \begin{pmatrix}
          t & a_0 \\
          -1 & a_1 + t
        \end{pmatrix}
        \Rightarrow \det(\tilde{A}) = t^2 + ta_1 + a_0
      \]
      and so clearly the determinant of $\tilde{A}$ is in the desired form. Now, assume that for any $\overline{A} \in \MM_{k \times k}(\FF)$ of the necessary form, $\det(\overline{A}) = t^{k} + t^{k-1}a_{k-1} + \cdots + t a_{1} + a_0$. We must now show that for $A \in \MM_{(k + 1) \times (k + 1)}(\FF)$, $\det(A) = t^{k+1} + t^{k}a_{k} + \cdots + t a_{1} + a_0$. However, 
      \begin{align*}
        \det(A) &= \begin{vmatrix}
          t & 0 & 0 & \cdots & 0 & a_0 \\
          -1 & t & 0 & \cdots & 0 & a_1 \\
          0 & -1 & t & \cdots & 0 & a_2 \\
          \vdots & \vdots & \vdots & & \vdots & \vdots \\
          0 & 0 & 0 & \cdots & -1 & a_{k} + t
        \end{vmatrix} \\
        &= t\begin{vmatrix}
          t & 0 & 0 & \cdots & 0 & a_1 \\
          -1 & t & 0 & \cdots & 0 & a_2 \\
          0 & -1 & t & \cdots & 0 & a_3 \\
          \vdots & \vdots & \vdots & & \vdots & \vdots \\
          0 & 0 & 0 & \cdots & -1 & a_{k} + t
        \end{vmatrix}
        + (-1)^{(k+1)+1}a_0\begin{vmatrix}
          -1 & t & 0 & \cdots & 0 \\
          0 & -1 & t & \cdots & 0 \\
          0 & 0 & -1 & \cdots & 0 \\
          \vdots & \vdots & \vdots & \ddots & 0 \\
          0 & 0 & 0 & \cdots & -1
        \end{vmatrix} \\
        &= t(t^{k} + t^{k-1}a_{k} + \cdots + t a_{2} + a_1) + (-1)^{k}(-1)^{k+2}a_0 \tag{IH with $a_1, \dots, a_k$; Problem~4.2.23} \\
        &= t^{k+1} + t^{k}a_k + \cdots + t^{2}a_2 + t a_1 + a_0
      \end{align*}
      which was to be shown. 
    \end{proof}
  \end{solution}

  \colorbox{red}{\underline{$4.4.2c$}:}
  \begin{solution}
    \[
      \begin{vmatrix}
        2 + i & -1 + 3i \\
        1 - 2i & 3 - i
      \end{vmatrix}
      = (2 + i)(3 - i) - (-1 +3i)(1 - 2i) = (7 + i) - (5 + 5i) = \boxed{2 - 4i}
    \]
  \end{solution}

  \newpage 

  \colorbox{red}{\underline{$4.4.3c$}:}
  \begin{solution}
    \[
      \begin{vmatrix}
        0 & 1 & 2 \\
        -1 & 0 & -3 \\
        2 & 3 & 0
      \end{vmatrix}
      = \begin{vmatrix}
        1 & 2 \\
        3 & 0
      \end{vmatrix}
      + 3\begin{vmatrix}
        0 & 1 \\
        2 & 3
      \end{vmatrix}
      = -6 + 3(-2) = \boxed{-12}
    \]
  \end{solution}

  \colorbox{red}{\underline{$4.4.3e$}:}
  \begin{solution}
    \[
      \begin{vmatrix}
        0 & 1 + i & 2 \\
        -2i & 0 & 1 - i \\
        3 & 4i & 0 
      \end{vmatrix}
      = 3\begin{vmatrix}
        1 + i & 2 \\
        0 & 1 - i
      \end{vmatrix}
      -4i\begin{vmatrix}
        0 & 2 \\
        -2i & 1 - i
      \end{vmatrix}
      = 3(2) - 4i(4i) = 6 + 16 = \boxed{22}
    \]
  \end{solution}

\end{document}